\documentclass{beamer}
\usepackage[utf8]{inputenc}

\usetheme{Boadilla}
\usecolortheme{lily}
\usepackage{amsmath,amssymb,amsfonts,amsthm}
\usepackage{mathtools}
\usepackage{txfonts}
\usepackage{tkz-euclide}
\usepackage{listings}
\usepackage{multicol}
\usepackage{adjustbox}
\usepackage{array}
\usepackage{tabularx}
\usepackage{lmodern}
\usepackage{gvv}
\usepackage{circuitikz}
\usepackage{tikz}
\usepackage{graphicx}

\setbeamertemplate{footline}
{
  \leavevmode%
  \hbox{%
  \begin{beamercolorbox}[wd=\paperwidth,ht=2.25ex,dp=1ex,right]{author in head/foot}%
    \insertframenumber{} / \inserttotalframenumber\hspace*{2ex} 
  \end{beamercolorbox}}%
  \vskip0pt%
}

\usepackage{tcolorbox}
\tcbuselibrary{minted,breakable,xparse,skins}




\providecommand{\nCr}[2]{\,^{#1}C_{#2}} % nCr
\providecommand{\nPr}[2]{\,^{#1}P_{#2}} % nPr
\providecommand{\mbf}{\mathbf}
\providecommand{\pr}[1]{\ensuremath{\Pr\left(#1\right)}}
\providecommand{\qfunc}[1]{\ensuremath{Q\left(#1\right)}}
\providecommand{\sbrak}[1]{\ensuremath{{}\left[#1\right]}}
\providecommand{\lsbrak}[1]{\ensuremath{{}\left[#1\right.}}
\providecommand{\rsbrak}[1]{\ensuremath{{}\left.#1\right]}}
\providecommand{\brak}[1]{\ensuremath{\left(#1\right)}}
\providecommand{\lbrak}[1]{\ensuremath{\left(#1\right.}}
\providecommand{\rbrak}[1]{\ensuremath{\left.#1\right)}}
\providecommand{\cbrak}[1]{\ensuremath{\left\{#1\right\}}}
\providecommand{\lcbrak}[1]{\ensuremath{\left\{#1\right.}}
\providecommand{\rcbrak}[1]{\ensuremath{\left.#1\right\}}}
\theoremstyle{remark}
\newcommand{\sgn}{\mathop{\mathrm{sgn}}}
\providecommand{\abs}[1]{\left\vert#1\right\vert}
\providecommand{\res}[1]{\Res\displaylimits_{#1}} 
\providecommand{\norm}[1]{\lVert#1\rVert}
\providecommand{\mtx}[1]{\mathbf{#1}}
\providecommand{\mean}[1]{E\left[ #1 \right]}
\providecommand{\fourier}{\overset{\mathcal{F}}{ \rightleftharpoons}}
%\providecommand{\hilbert}{\overset{\mathcal{H}}{ \rightleftharpoons}}
\providecommand{\system}{\overset{\mathcal{H}}{ \longleftrightarrow}}
	%\newcommand{\solution}[2]{\textbf{Solution:}{#1}}
%\newcommand{\solution}{\noindent \textbf{Solution: }}
\providecommand{\dec}[2]{\ensuremath{\overset{#1}{\underset{#2}{\gtrless}}}}
\newcommand{\myvec}[1]{\ensuremath{\begin{pmatrix}#1\end{pmatrix}}}
\let\vec\mathbf

\lstset{
%language=C,
frame=single, 
breaklines=true,
columns=fullflexible
}

\numberwithin{equation}{section}

\lstset{
  language=Python,
  basicstyle=\ttfamily\small,
  keywordstyle=\color{blue},
  stringstyle=\color{orange},
  numbers=left,
  numberstyle=\tiny\color{gray},
  breaklines=true,
  showstringspaces=false
}

\numberwithin{equation}{section}

\title{12.70}
\author{Rushil Shanmukha Srinivas \\EE25BTECH11057 \\ Electrical Enggineering ,\\IIT Hyderabad.}

\date{\today} 
\begin{document} 

\begin{frame}
\titlepage
\end{frame}

\section*{Outline}
\begin{frame}
\tableofcontents
\end{frame}
\section{Problem}
\begin{frame}
\frametitle{Problem Statement}
\textbf{Question :} 
$X$ bullocks and $Y$ tractors take 8 days to plough a field. If we halve the number of bullocks and double the number of tractors, it takes 5 days to plough the same field.
How many days will it take $X$ bullocks alone to plough the field?

\end{frame}
\section{Solution}
\subsection{Number of days and Work}
\begin{frame}
\frametitle{Number of days and Work}
%\framesubtitle{Literature}
The total work done = W

Work done by 1 bullock in one day = x

Work done by 1 tractor in one day = y

Given,
\begin{align}
\myvec{8X&8Y}\myvec{x\\y} = W
\end{align}
\begin{align}
\myvec{\frac{5X}{2}&10Y}\myvec{x\\y} = W
\end{align}
The combined matrix is
\begin{align}
\myvec{8X&8Y\\
        \frac{5X}{2}&10Y}\myvec{x\\y} = \myvec{W\\W}
\end{align}
Forming the augmented matrix ,
\begin{align}
  \myaugvec{2}{8X & 8Y & W\\ \frac{5X}{2} & 10Y & W}
\end{align}
\end{frame}
\begin{frame}
Using Gaussian Eliminations
\begin{align}
\myaugvec{2}{8X & 8Y & W\\ \frac{5X}{2} & 10Y & W} 
\xleftrightarrow{\;R_2 \to 4R_2 -5R_1}
 \myaugvec{2}{8X & 8Y & W\\-30X & 0 & -W} 
\end{align}
\begin{align}
30Xx = W
\end{align}
Number of days taken = 30

\end{frame}
\section{C Code}
\begin{frame}[fragile]
\frametitle{C Code }
\begin{lstlisting}[language=C]
#include <stdio.h>

double days_for_bullocks_alone(double X, double Y) {
    // From equations:
    // 8*(X*b + Y*t) = 1
    // 5*((X/2)*b + 2*Y*t) = 1
    // Solve for b, t symbolically and return 1/(X*b)
    double a1 = 8 * X;
    double b1 = 8 * Y;
    double c1 = 1;

    double a2 = 5 * (X / 2.0);
    double b2 = 5 * (2.0 * Y);
    double c2 = 1;

    // Expressing equations in rate form:
    // (a1/a1)*b + (b1/a1)*t = c1/a1
    // (a2/a1)*b + (b2/a1)*t = c2/a1
\end{lstlisting}
\end{frame}
\begin{frame}[fragile]
\begin{lstlisting}
    double det = a1 * b2 - a2 * b1;
    double b_rate = (c1 * b2 - c2 * b1) / det;
    double t_rate = (a1 * c2 - a2 * c1) / det;

    // time for X bullocks alone = 1 / (X * b_rate)
    double T = 1.0 / (X * b_rate);
    return T;
}
int main() {
    double X, Y;
    printf("Enter number of bullocks (X): ");
    scanf("%lf", &X);
    printf("Enter number of tractors (Y): ");
    scanf("%lf", &Y);

    double T = days_for_bullocks_alone(X, Y);
    printf("Days taken by %g bullocks alone = %.2f days\n", X, T);

    return 0;
}

\end{lstlisting}
\end{frame}

\section{Python Code}
\begin{frame}[fragile]
\frametitle{call.py}
\begin{lstlisting}[language=Python]
import ctypes

# Load the shared object file (make sure plough.so is in the same directory)
plough = ctypes.CDLL('./plough.so')

# Define argument and return types for the C function
plough.days_for_bullocks_alone.argtypes = [ctypes.c_double, ctypes.c_double]
plough.days_for_bullocks_alone.restype = ctypes.c_double

def main():
    print(" Bullocks and Tractors Ploughing Problem ")
    # Get input values
    X = float(input("Enter number of bullocks (X): "))
    Y = float(input("Enter number of tractors (Y): "))

    # Call the C function
    T = plough.days_for_bullocks_alone(X, Y)
\end{lstlisting}
\end{frame}
\begin{frame}[fragile]
\begin{lstlisting}
    # Display result
    print(f"\n➡ Days taken by {X:.0f} bullocks alone to plough the field = {T:.2f} days")

if __name__ == "__main__":
    main()
\end{lstlisting}
\end{frame}

\end{document}

